%% ---------------------------------------------------------------------------
%% paNotation.tex
%%
%% Notation
%%
%% $Id: paNotation.tex,v 1.15 2004/03/30 05:55:59 alvarado Exp $
%% ---------------------------------------------------------------------------

\cleardoublepage
\renewcommand{\nomname}{Lista de símbolos y abreviaciones}
\markboth{\nomname}{\nomname}
\renewcommand{\nompreamble}{\addcontentsline{toc}{chapter}{\nomname}%
\setlength{\nomitemsep}{-\parsep}
\setlength{\itemsep}{10ex}
}

%%
% Algunas abreviaciones
%%

\syma{TP}{Test Program}
\syma{PPV}{Product Platform Validation}
\syma{DTD}{Debug and Test Development}
\syma{DTD}{Debug and Test Development}
\syma{ALU}{Unidad Aritmético-Lógica}
\syma{AVX-512}{Advanced Vector Extensions 512}
\syma{L1}{Caché de Nivel 1}
\syma{L2}{Caché de Nivel 2}
\syma{L3}{Caché de Nivel 3}
\syma{QPI}{QuickPath Interconnect}
\syma{UPI}{Ultra Path Interconnect}
\syma{VT-d}{Virtualization Technology for Directed I/O}
\syma{VT-x}{Virtualization Technology}
\syma{CPU}{Central Processing Unit} 
\syma{Vdd}{Voltage Drain Drain} 
\syma{GHz}{Gigahertz} 
\syma{MHz}{Megahertz} 
\syma{MOSFET}{Metal-Oxide-Semiconductor Field-Effect Transistor} 
\syma{CMOS}{Complementary Metal-Oxide-Semiconductor} 
\syma{VLSI}{Very-Large-Scale Integration} 
\syma{ASIC}{Application-Specific Integrated Circuit} 
\syma{VTC}{Voltage Transfer Characteristic} 
\syma{AC}{Alternating Current} 
\syma{DC}{Direct Current} 
\syma{SoC}{System on Chip} 
\syma{PVT}{Process, Voltage, Temperature} 
\syma{TDP}{Thermal Design Power} 
\syma{HPC}{High-Performance Computing} 
\syma{EDA}{Electronic Design Automation} 
\syma{RTL}{Register-Transfer Level}
\printnomenclature[20mm]

%%% Local Variables:
%%% mode: latex
%%% TeX-master: "paMain"
%%% End:
